\subsection{三角関数型}
またもや似た見た目の漸化式を考える.
\[
  a_1=\frac{\sqrt{3}}{2},\quad a_{n+1}=2a_n^2-1
\]

この漸化式も前節と同様に,今までの方法ではうまく解くことができない.

今回の鍵は三角関数である.
三角関数と置くことで,数列の漸化式が角度の漸化式へ変わることがある.
その仕組みを確認するために,まず加法定理について復習しよう.
\begin{align*}
  \sin(\alpha+\beta)
    &=\sin\alpha\cos\beta+\cos\alpha\sin\beta\\
  \cos(\alpha+\beta)
    &=\cos\alpha\cos\beta-\sin\alpha\sin\beta\\
  \tan(\alpha+\beta)
    &=\frac{\tan\alpha+\tan\beta}
            {1-\tan\alpha\tan\beta}
\end{align*}
ここで $\alpha=\beta=x$ とすれば,倍角公式
\begin{align*}
  \sin 2x&=2\sin x\cos x\\
  \cos 2x&=2\cos^2x-1=1-2\sin^2x\\
  \tan 2x&=\frac{2\tan x}{1-\tan^2x}
\end{align*}
を得る.また,倍角公式から
\begin{align*}
  \sin^2x&=\frac{1-\cos 2x}{2}\\
  \cos^2x&=\frac{1+\cos 2x}{2}\\
  \tan^2x&=\frac{1-\cos 2x}{1+\cos 2x}
\end{align*}
を導くことができる.また,加法定理から,\,$3x=2x+1$とすることで,
\begin{align*}
  \sin 3x&=3\sin x-4\sin^3x\\
  \cos 3x&=4\cos^3x-3\cos x
\end{align*}
が求められる.

これらを踏まえて次の漸化式を考える.
\[
  a_1=\frac{\sqrt{3}}{2},\quad a_{n+1}=2a_n^2-1.
\]
右辺は $\cos 2x=2\cos^2x-1$ と一致している.さらに,
初項は
\[
  a_1=\cos\frac{\pi}{6}
\]
と書ける.そこで $a_n=\cos\theta_n$ と置くと,
\[
  \theta_1=\frac{\pi}{6},
  \quad
  \theta_{n+1}=2\theta_n
\]
と考えられる.従って候補は
\[
  a_n=\cos\left(2^{n-1}\frac{\pi}{6}\right)
\]
である.

この候補を帰納法で確認する.\,$n=1$ では初項に一致する.ある $n$ で
\[
  a_n=\cos\left(2^{n-1}\frac{\pi}{6}\right)
\]
が成り立つとすると,
\begin{align*}
  a_{n+1}
    &=2a_n^2-1\\
    &=2\cos^2\left(2^{n-1}\frac{\pi}{6}\right)-1\\
    &=\cos\left(2^n\frac{\pi}{6}\right).
\end{align*}
これは $n+1$ の式である.よって,
\[
  a_n=\cos\left(2^{n-1}\frac{\pi}{6}\right)
\]
が示された.

この解法の本質は,二次式を無理に展開したり因数分解したりすることではない.
数列の値を角度の関数として表すことで,二次式の計算を角度の2倍という単純な操作に翻訳しているのである.

今度は
\[
  a_1=\tan\frac{\pi}{12}=2-\sqrt{3},
  \quad
  a_{n+1}=\frac{2a_n}{1-a_n^2}
\]
を考える.右辺は$\tan$の倍角公式そのものである.したがって,
\[
  a_n=\tan\left(2^{n-1}\frac{\pi}{12}\right)
\]
と予想できる.

実際,\,$n=1$ で初項に一致し,帰納法の段階では
\begin{align*}
  a_{n+1}
    &=\frac{2a_n}{1-a_n^2}\\
    &=\cfrac{2\tan\left(2^{n-1}\cfrac{\pi}{12}\right)}
            {1-\tan^2\left(2^{n-1}\cfrac{\pi}{12}\right)}\\
    &=\tan\left(2^n\frac{\pi}{12}\right)
\end{align*}
となる.ただし,分母が $0$ になると$\tan$は定義できない.
三角関数型では,置換を思いつくだけでなく,途中で定義できなくなる項がないかも確認しなければならない.

余弦の例と$\tan$の例は,同じ方針を共有している.すなわち,
\[
  a_n=\varphi(\theta_n),
  \quad
  \varphi(m\theta)=F_m(\varphi(\theta))
\]
となる関数 $\varphi$ と整数倍公式を探し,角度の漸化式を解いている.
いわば媒介変数表示や置換のようなことをしている.

また,ここで現れた多項式
\[
  2x^2-1,
  \quad
  4x^3-3x
\]
は偶然の産物ではない.チェビシェフ多項式 $T_m(x)$ を
\[
  T_m(\cos\theta)=\cos(m\theta)
\]
で定義すれば,
\[
  T_0(x)=1,
  \quad
  T_1(x)=x,
  \quad
  T_{m+1}(x)=2xT_m(x)-T_{m-1}(x)
\]
となる.例えば $T_2(x)=2x^2-1$ である.
したがって,余弦の例は
\[
  a_{n+1}=T_2(a_n),
  \quad
  a_n=T_{2^{n-1}}(a_1)
\]
と書き直せる.三角関数の公式から始めて多項式の族へ進む道が,
付録のチェビシェフ多項式へつながっている.
